\documentclass{article}
\usepackage{lmodern}
\usepackage{amsmath}
\usepackage{listings}
\usepackage{fullpage}
\usepackage{parskip}
\usepackage{xcolor}
\lstset{
	basicstyle=\ttfamily,
	columns=fullflexible,
	frame=single,
	breaklines=true,
	postbreak=\mbox{\textcolor{red}{$\hookrightarrow$}\space},
}
\begin{document}
\title{Experiment `p\_4\_2\_a\_min\_max\_array' Results}
\author{\today}
\date{}
\maketitle
\textbf{Experiment outcome: }SUCCESS
\\\textbf{Bad responses: }0
\\\textbf{Responses containing}~\texttt{assume}~\textbf{: }0
\\\textbf{Responses containing}~\texttt{expect}~\textbf{: }0
\\\textbf{Resolution attempts: }1
\\\textbf{Hard fails (resolution): }0
\\\textbf{Soft fails (resolution): }0
\\\textbf{Verification attempts: }1
\section*{Problem Specification}
\textbf{Problem name: }p\_4\_2\_a\_min\_max\_array
\\\textbf{Natural language statement: }null
\\\textbf{Method signature: }p\_4\_2\_a\_min\_max\_array(inputs: array<int>) returns (min: int, max: int)
\subsection*{Ensures}
\begin{itemize}
\item \texttt{forall i :: 0 <= i < inputs.Length ==> min <= inputs[i]}
\item \texttt{forall i :: 0 <= i < inputs.Length ==> max >= inputs[i]}
\item \texttt{exists i :: 0 <= i < inputs.Length \&\& min == inputs[i]}
\item \texttt{exists i :: 0 <= i < inputs.Length \&\& max == inputs[i]}
\end{itemize}
\subsection*{Requires}
\begin{itemize}
\item \texttt{0 < inputs.Length}
\end{itemize}
\clearpage
\section*{GenAI interactions}
Below you will find all interactions between the `user' (program) and the `assistant' (OpenAI).
\subsection*{Program $\rightarrow$ GenAI}
\begin{lstlisting}
You are given the following task to perform in Dafny:

Write a method to satisfy the below requirements.

The signature should be:

p_4_2_a_min_max_array(inputs: array<int>) returns (min: int, max: int)

The method should respect the following contract:

requires 0 < inputs.Length, ensures forall i :: 0 <= i < inputs.Length ==> min <= inputs[i], ensures forall i :: 0 <= i < inputs.Length ==> max >= inputs[i], ensures exists i :: 0 <= i < inputs.Length && min == inputs[i], ensures exists i :: 0 <= i < inputs.Length && max == inputs[i]

Produce and show only the Dafny body of this method, including the curly braces that surround it. Do not show the signature nor contract. You must not use 'assume' nor 'expect' anywhere in your code.
\end{lstlisting}
\subsection*{GenAI $\rightarrow$ Program}
\textbf{System fingerprint: }null
\\\textbf{ID: }chatcmpl-DnLI7toHsksmHlKuFbylkLNi2MKWj
\\\textbf{Model: }gpt-5.5-2026-04-23
\\\textbf{Created at: }1780652287
\\\textbf{Finish reason: }stop
\\\textbf{Usage: }392 tokens in, and 736 tokens out
\begin{lstlisting}
{
  min := inputs[0];
  max := inputs[0];
  var idxMin := 0;
  var idxMax := 0;
  var i := 1;
  while i < inputs.Length
    invariant 1 <= i <= inputs.Length
    invariant 0 <= idxMin < i
    invariant 0 <= idxMax < i
    invariant min == inputs[idxMin]
    invariant max == inputs[idxMax]
    invariant forall j :: 0 <= j < i ==> min <= inputs[j]
    invariant forall j :: 0 <= j < i ==> max >= inputs[j]
  {
    if inputs[i] < min {
      min := inputs[i];
      idxMin := i;
    }
    if inputs[i] > max {
      max := inputs[i];
      idxMax := i;
    }
    i := i + 1;
  }
}
\end{lstlisting}
\section*{Final Program}
\begin{lstlisting}
method p_4_2_a_min_max_array(inputs: array<int>) returns (min: int, max: int)
	requires 0 < inputs.Length
	ensures forall i :: 0 <= i < inputs.Length ==> min <= inputs[i]
	ensures forall i :: 0 <= i < inputs.Length ==> max >= inputs[i]
	ensures exists i :: 0 <= i < inputs.Length && min == inputs[i]
	ensures exists i :: 0 <= i < inputs.Length && max == inputs[i]
{
  min := inputs[0];
  max := inputs[0];
  var idxMin := 0;
  var idxMax := 0;
  var i := 1;
  while i < inputs.Length
    invariant 1 <= i <= inputs.Length
    invariant 0 <= idxMin < i
    invariant 0 <= idxMax < i
    invariant min == inputs[idxMin]
    invariant max == inputs[idxMax]
    invariant forall j :: 0 <= j < i ==> min <= inputs[j]
    invariant forall j :: 0 <= j < i ==> max >= inputs[j]
  {
    if inputs[i] < min {
      min := inputs[i];
      idxMin := i;
    }
    if inputs[i] > max {
      max := inputs[i];
      idxMax := i;
    }
    i := i + 1;
  }
}
\end{lstlisting}
\section*{Total Token Usage}
\textbf{Input tokens: }392
\\\textbf{Output tokens: }736
\\\textbf{Reasoning tokens: }512
\\\textbf{Sum of `total tokens': }1128
\section*{Experiment Timings}
\textbf{Overall Experiment} started at 1780652287040, ended at 1780652312103, lasting 25063ms (25.06 seconds)
\\\textbf{Iteration \#1} started at 1780652287041, ended at 1780652312103, lasting 25062ms (25.06 seconds)
\end{document}
